\section{Toric tangent bundles}
\label{sec:toric}

The geometric inputs in this section are proved in the companion
\cite{Chambers}. We state them explicitly to distinguish the geometry
needed for Theorem~\ref{thm:main} from the analytic conclusion.

For $k\ge1$ let
\[
 X_k=\PP_{(\PP^1)^{2k-1}}
       \bigl(\mathcal O\oplus\mathcal O(1^k,-1^{k-1})\bigr),
\]
using quotient projectivization. The twisting line bundle has degree
$1$ on the first $k$ base factors and degree $-1$ on the remaining
$k-1$. Let $\pi:X_k\to(\PP^1)^k$ be the projection onto the
first $k$ base factors, with components $\pi_j$.
This is a smooth toric Fano variety of dimension and Picard number
$2k$. Put $D=-K_{X_k}$ and $h_j=c_1(\pi_j^*\mathcal O_{\PP^1}(1))$.
The tangent bundle fits into
\begin{equation}\label{eq:toric-extension}
 0\longrightarrow F_0:=T_{X_k/(\PP^1)^k}
 \longrightarrow T_{X_k}
 \longrightarrow\bigoplus_{j=1}^k\pi_j^*\mathcal O_{\PP^1}(2)
 \longrightarrow0.
\end{equation}
At $D$, $F_0$ is stable of rank $k$, the quotient line bundles are
stable of the same slope, and all factors are pairwise nonisomorphic.
Every component of the extension class is nonzero. The functions
\[
 z_j(H)=\left(2h_j-\frac{D}{2k}\right)H^{2k-1}
\]
extend to independent real analytic local coordinates at $D$.
The local stable locus is $z_j>0$ for every $j$.

Choose smooth K\"ahler representatives of the nearby ample classes.
The slopes in \cite{Chambers} are intersection numbers without a
factorial denominator. In the present normalization,
\begin{equation}\label{eq:toric-epsilon}
 \epsilon_j(H)=\frac{2\pi z_j(H)}{(2k-1)!V_H}
              =\frac{4\pi k\,z_j(H)}{H^{2k}}.
\end{equation}
At the wall the differential of this map is a common positive
multiple of $dz_j$, so the $\epsilon_j$ are also independent local
coordinates after completion by $k$ further coordinates.

\begin{corollary}\label{cor:toric}
For the HE tangent bundle of $X_k$ in its nearby stable chamber,
there are exactly $k$ small trace-free Dolbeault eigenvalues. To
relative error $O(\sum_j\epsilon_j)$ they are the eigenvalues of
\[
 \diag(\epsilon_1,\ldots,\epsilon_k)+\frac1k vv^*,
 \qquad v_j=\sqrt{\epsilon_j}.
\]
Along a path with every $\epsilon_j=\epsilon$, the leading spectrum
is $\epsilon$ with multiplicity $k-1$ and $2\epsilon$ once.
Positive paths with arbitrary distinct orders of vanishing of the
$\epsilon_j$ are realized by genuine nearby polarizations.
\end{corollary}

\begin{proof}
Apply Theorem~\ref{thm:main} with $r_0=k$ and $r_j=1$.
For equal surpluses the rank-one perturbation has eigenvalue
$\epsilon$ on the complement of $(1,\ldots,1)$ and $2\epsilon$
on its span. Independent local coordinates realize the indicated
positive paths while staying in the ample cone.
\end{proof}

\begin{example}
For $X_2$, write the positive surpluses as $\epsilon_1,\epsilon_2$.
The two comparison eigenvalues are
\[
 \nu_\pm=\frac34(\epsilon_1+\epsilon_2)
    \pm\frac14\sqrt{9(\epsilon_1-\epsilon_2)^2+4\epsilon_1\epsilon_2}.
\]
If $\epsilon_1/\epsilon_2\to0$, the actual HE eigenvalues satisfy
$\lambda_1\sim\frac43\epsilon_1$ and
$\lambda_2\sim\frac32\epsilon_2$. The relative theorem applies
equally when $\epsilon_1=e^{-1/\epsilon_2^2}$.
\end{example}

\begin{example}
For $X_3$ choose $\epsilon_1=s^6$, $\epsilon_2=s^4$ and
$\epsilon_3=s^2$, $s>0$. The characteristic polynomial of the
comparison matrix is
\[
 \lambda^3-\frac43(\epsilon_1+\epsilon_2+\epsilon_3)\lambda^2
 +\frac53(\epsilon_1\epsilon_2+\epsilon_1\epsilon_3+
                         \epsilon_2\epsilon_3)\lambda
 -2\epsilon_1\epsilon_2\epsilon_3.
\]
Its largest root is asymptotic to $\frac43s^2$. The second elementary
symmetric function then gives the middle root $\frac54s^4$, and
the determinant gives the smallest root $\frac65s^6$.
Consequently the actual tangent-bundle spectrum satisfies
\[
 \lambda_1\sim\tfrac65s^6,\qquad
 \lambda_2\sim\tfrac54s^4,\qquad
 \lambda_3\sim\tfrac43s^2.
\]
These are changes of polarization on one fixed six-dimensional variety.
\end{example}

Theorem~4.1 and equation~(16) of \cite{Chambers} supply further examples
from the strictly semistable smooth toric Fano census in dimensions
three through six, overlapping this family at $X_2$ and $X_3$.
All 31 indecomposable cases give star extensions of distinct locally
free stable factors, with nonzero components. Twenty-six have two
factors and hence also
satisfy Theorem~\ref{thm:second-order}. The five cases with more
factors have ranks
\[
 (2,1,1),\ (3,1,1),\ (4,1,1),\ (4,1,1),\ (3,1,1,1).
\]
In dimension $n$, for a quotient of rank $r_j$ with first Chern class
$C_j$, the general conversion is
\[
 \epsilon_j(H)=\frac{2\pi}{r_j(n-1)!V_H}
       \left(C_j-\frac{r_j}{n}c_1(T_X)\right)H^{n-1}.
\]
Computing the two-factor
coefficient $D$ would additionally require the wall Green operator
and normalized harmonic extension; the toric degree data do not
provide those analytic quantities.

\section{Scope and further questions}

The star structure is responsible for the disappearance of harmonic
norms from the leading spectrum. It gives one positive scaling
variable for each independent scalar equation, and a split summand
on every coordinate face. For general two-step extension graphs
with cycles, the number of edge weights can exceed the number of
independent moment equations; the ranks and slope surpluses then
need not determine all leading weights. Repeated stable factors
also replace block scalars by noncommutative endomorphism algebras.
Neither setting is included in Theorem~\ref{thm:main}.

The estimates concern equilibrium metrics for varying K\"ahler forms.
Relating them to the long-time HYM flow as the polarization also
approaches a wall requires a uniform two-parameter analysis. The
inverse-time behavior on the wall and the positive equilibrium gap
inside the chamber do not by themselves justify interchanging those
limits.
